% https://www.latex-project.org/help/documentation/
% https://www.learnlatex.org/

% preamble
\documentclass[twocolumn]{article}
\usepackage[T1]{fontenc}

% \usepackage[french]{babel}
% \newcommand\kw[1]{\textbf{\itshape #1}}


\usepackage{xcolor}

\newcommand\cred[1]{\textcolor{red}{\itshape #1}}
\newcommand\cblu[1]{\textcolor{blue}{\itshape #1}}


\usepackage{lipsum}


\usepackage[margin = 1in]{geometry} % To force hyphenation here


\begin{document}
% There is a titlepage environment for when you want to do custom design, 
% but this is out of the scope of this introduction. 
% If you want to do your own document designs you can either 
% use a customisable class, such as memoir, or start with 
% one of LaTeX’s base classes, like book and use it as a starting point.
\author{sol4nki \and yeehaw}
\title{ASSIGNMENT-2}
\date{18st August 2026}
\maketitle
so hello there!

Something about \cred{apples} and \cblu{oranges}.

\tableofcontents


\lipsum
\lipsum
\lipsum
\\
\\
\\
This is a lot of filler which is going to demonstrate how LaTeX hyphenates
material, and which will be able to give us at least one hyphenation point.
This is a lot of filler which is going to demonstrate how LaTeX hyphenates
material, and which will be able to give us at least one hyphenation point.


testing lorem ipsum
\section{First section}
Text of material in the first section

second para.

\subsection{subsection of 1st section}
hello there!
\paragraph{hello world}
\subparagraph{yes?}
2nd para for no reason

\section{random stuff} 

Some text with \emph{emphasis and \emph{nested} content}.

Some text in \textit{italic and \textit{nested} content}.

Some text in \textbf{bold and \textbf{nested} content}.

\section{lists}

\subsection{Ordered}
\begin{enumerate}
  \item An entry
  \item Another One
  \item Wow! Three entries
\end{enumerate}

\subsection{Unordered}
\begin{itemize}
  \item An entry
  \item Another One
  \item Wow! Three entries
\end{itemize}

\subsection{Descriptive List}
\begin{description}
\item[Dog:] member of the genus Canis, which forms part of the wolf-like canids,
  and is the most widely abundant terrestrial carnivore.
\item[Cat:] domestic species of small carnivorous mammal. It is the only
  domesticated species in the family Felidae and is often referred to as the
  domestic cat to distinguish it from the wild members of the family.
\end{description}


\end{document}


